\chapter{1960:Willard F. Libby}

\label{chap:chemistry1960}

The Nobel Prize in Chemistry for the year 1960 was awarded to Willard F. Libby.

\textbf{Motivation:}

for his method to use carbon-14 for age determination in archaeology, geology, geophysics, and other branches of science

\section{Willard F. Libby}

\subsection*{Early Life and Education}

Willard F. Libby was born on 1908-12-17 in Grand Valley, Colorado, USA.

\subsection*{Career and Major Research}

Libby studied at the University of California, Berkeley, where he earned his doctorate in 1933, and then taught there. During the Second World War he worked on the Manhattan Project in the isotopic separation of uranium. After the war, at the University of Chicago, he developed the method of radiocarbon dating, announcing in 1949 the first determinations of the age of archaeological samples based on the decay of carbon-14.

\subsection*{Nobel Prize-winning Work}

The work for which Willard F. Libby was awarded the Nobel Prize in Chemistry in 1960 was described by the Nobel Committee as: "for his method to use carbon-14 for age determination in archaeology, geology, geophysics, and other branches of science".

Libby realized that carbon-14, produced continuously in the atmosphere by cosmic radiation and incorporated into living organisms through photosynthesis, decays with a known half-life after death. By measuring the residual carbon-14 in ancient wood, charcoal, and bone, he could calculate the time elapsed since the material stopped exchanging carbon with the atmosphere.

\subsection*{Significance and Impact}

Radiocarbon dating provided, for the first time, a general method of absolute dating for organic remains up to about fifty thousand years old. It revolutionized archaeology, geology, and Quaternary science, and its checks against objects of known age proved the method reliable.

\subsection*{Later Career and Honors}

Libby left Chicago to become a member of the United States Atomic Energy Commission in 1955 and later was a professor of chemistry at the University of California, Los Angeles. He died on 1980-09-08 in Los Angeles.

\subsection*{Legacy}

Radiocarbon dating remains one of the most important dating methods in archaeology and earth science, and Libby's discovery founded the field of radiometric dating with cosmogenic isotopes.

\subsection*{Modern Developments}

Libby's radiocarbon dating revolutionized archaeology, geology, and paleoclimatology, providing a universal clock for organic materials up to about 50,000 years old. Radiocarbon dating is now routine in dating ancient artifacts, studying climate change, and forensic science, and Libby's methods extended to other cosmogenic isotopes.

\subsection*{Selected Publications}

\begin{itemize}

\item \emph{Radiocarbon Dating}, University of Chicago Press, 1952.

\end{itemize}

\clearpage

\section*{Chapter Summary}

The 1960 prize recognized the development of radiocarbon dating. Willard F. Libby showed at the University of Chicago that the decay of carbon-14 in organic material could be used to determine its age, providing archaeologists and geologists with a method of absolute dating for organic remains.

